\chapter{Hermitian Perfect Objects and Homotopy-Symmetric Poincar\'e Structures}
\label{chap:hermitian-poincare}

\section{Purpose and standing conventions}
\label{sec:hermitian-poincare-purpose}

This chapter begins with the rigid stable symmetric monoidal categories and
coherent dualities constructed in
\Ref{chap:perfect-rigid-enrichment}. Thus, for a commutative algebra object
\[
    A\in\CAlg(\V),
\]
and for a relative presheaf
\[
    X\in\PSh(\Aff),
\]
we have coherent perfect dualities
\[
    D_A=(-)^\vee
    \colon
    \Perf(A)^{\op}
    \xrightarrow{\;\simeq\;}
    \Perf(A)
\]
and
\[
    D_X=(-)^\vee
    \colon
    \Perf(X)^{\op}
    \xrightarrow{\;\simeq\;}
    \Perf(X).
\]

The purpose of the chapter is to study coherent homotopy-fixed points of
these dualities and to identify them with the nondegenerate objects of the
homotopy-symmetric Poincar\'e functors
\[
    \mathcal Q_A^{\mathrm s}(M)
    :=
    \mathbb B_A(M,M)^{hC_2}
\]
and
\[
    \mathcal Q_X^{\mathrm s}(M)
    :=
    \mathbb B_X(M,M)^{hC_2}.
\]
The represented bilinear functors are
\[
    \mathbb B_A(M,N)
    \simeq
    \map_{\Perf(A)}(M,N^\vee)
\]
and
\[
    \mathbb B_X(M,N)
    \simeq
    \map_{\Perf(X)}(M,N^\vee).
\]

No involution on
\[
    \V,
    \qquad
    A,
    \qquad\text{or}\qquad
    \Aff_{\V}
\]
is assumed. A Hermitian structure in this chapter means a coherent
homotopy-fixed-point structure for the canonical rigid duality whose
underlying map is an equivalence
\[
    M
    \xrightarrow{\;\simeq\;}
    M^\vee.
\]
It satisfies biduality symmetry together with all higher
\[
    C_2
\]
-coherences. No semilinear conjugation or positivity condition is included.

The chapter proceeds from coherent fixed points and their induced dagger,
through cross-effects and polarization, to represented Poincar\'e
structures. It then proves that the homotopy-symmetric refinement of any
coherent perfect duality has Poincar\'e objects precisely its coherent fixed
points, applies this theorem to perfect modules and perfect objects on
relative stacks, and establishes pullback and base-change functoriality.

The categorical, mapping-space, coherence, universe, and sheaf conventions
of
\Ref{chap:relative-ag-monoidal-syntax}
and
\Ref{chap:perfect-rigid-enrichment}
remain in force.

\begin{convention}
\label{conv:equivalence-mapping-space}
For objects
\[
    X,Y\in\mathcal C,
\]
write
\[
    \Map^\simeq_{\mathcal C}(X,Y)
    \subseteq
    \Map_{\mathcal C}(X,Y)
\]
for the union of those connected components represented by equivalences.
Equivalently,
\[
    \Map^\simeq_{\mathcal C}(X,Y)
    =
    \Map_{\mathcal C^\simeq}(X,Y),
\]
where
\[
    \mathcal C^\simeq
\]
is the maximal $\infty$-groupoid of
\[
    \mathcal C.
\]
\end{convention}


\begin{convention}[Higher coherence in the Hermitian and Poincar\'e layers]
\label{conv:rigid-poincare-higher-coherence}
In addition to
\Ref{conv:higher-coherence},
involutivity, cross-effect symmetry, polarization, homotopy-fixed-point,
and Poincar\'e data are always understood together with all required
coherent higher homotopies.

Whenever such data are suppressed from the notation, their existence is
supplied by a symmetric-bilinear, homotopy-fixed-point, or
Goodwillie-excisive construction, by a universal property with contractible
choice space, or by an explicitly cited coherence theorem.
\end{convention}

\begin{convention}[Double-duality notation]
\label{conv:double-duality-notation}
Let
\[
    D
    \colon
    \mathcal C^{\op}
    \longrightarrow
    \mathcal C
\]
be a contravariant functor, and let
\[
    D^{\op}
    \colon
    \mathcal C
    \longrightarrow
    \mathcal C^{\op}
\]
be its opposite functor.

We use
\[
    D^2
\]
as an abbreviation for the correctly typed covariant composite
\[
    D D^{\op}
    \colon
    \mathcal C
    \longrightarrow
    \mathcal C.
\]
Thus
\[
    D^2X
    :=
    D D^{\op}(X).
\]

More generally, if
\[
    D_{\mathcal Q}
    \colon
    \mathcal C^{\op}
    \longrightarrow
    \mathcal C
\]
is the functor representing the cross-effect of a quadratic functor, then
\[
    D_{\mathcal Q}^2
    :=
    D_{\mathcal Q}D_{\mathcal Q}^{\op}.
\]

All occurrences of iterated contravariant duality in this chapter are
interpreted using this convention.
\end{convention}

\section{Hermitian perfect modules and the induced dagger}
\label{sec:hermitian-perfect-modules}

The coherent perfect duality of
\Ref{prop:canonical-perfect-duality}
induces a covariant
\[
    C_2
\]
-action on the maximal $\infty$-groupoid
\[
    \Perf(A)^\simeq.
\]
It sends an object
\[
    M
\]
to
\[
    M^\vee,
\]
and an equivalence
\[
    f
    \colon
    M
    \xrightarrow{\;\simeq\;}
    N
\]
to
\[
    (f^{-1})^\vee
    \colon
    M^\vee
    \xrightarrow{\;\simeq\;}
    N^\vee.
\]


\subsection{Hermitian perfect modules}

\begin{definition}[Hermitian perfect module]
\label{def:hermitian-perfect-module}
A \emph{nondegenerate Hermitian structure} on
\[
    M\in\Perf(A)
\]
is a coherent homotopy-fixed-point structure for the duality action on
\[
    \Perf(A)^\simeq.
\]

Its underlying equivalence is a map
\[
    h_M
    \colon
    M
    \xrightarrow{\;\simeq\;}
    M^\vee
\]
whose first fixed-point coherence is
\[
    h_M^\vee
    \circ
    \operatorname{bid}_M
    \simeq
    h_M.
\]
A \emph{Hermitian perfect module} is a pair
\[
    (M,h_M)
\]
equipped with this coherent fixed-point structure.
\end{definition}

\begin{remark}
\label{rem:hermitian-no-conjugation}
The adjective Hermitian refers here to coherent symmetric self-duality in
the rigid category
\[
    \Perf(A).
\]
Over a field with no chosen conjugation, this produces symmetric bilinear
rather than sesquilinear forms.

The terminology of this chapter reserves the phrase
\emph{Hermitian perfect module} for a nondegenerate coherently symmetric
self-duality. In the broader Poincar\'e-$\infty$-categorical terminology,
an arbitrary point of a quadratic functor is also called a Hermitian form,
while the adjective Poincar\'e records nondegeneracy.
\end{remark}


\subsection{The induced dagger}

\begin{proposition}[Adjoint of a morphism]
\label{prop:dagger-on-perfect-modules}
Let
\[
    (M,h_M)
    \qquad\text{and}\qquad
    (N,h_N)
\]
be Hermitian perfect modules. For
\[
    f
    \colon
    M
    \longrightarrow
    N,
\]
define
\[
    f^\dagger
    :=
    h_M^{-1}
    \circ
    f^\vee
    \circ
    h_N
    \colon
    N
    \longrightarrow
    M.
\]
Then
\[
    (g\circ f)^\dagger
    \simeq
    f^\dagger\circ g^\dagger,
\]
\[
    (f^\dagger)^\dagger
    \simeq
    f,
\]
and
\[
    \id_M^\dagger
    \simeq
    \id_M.
\]
\end{proposition}

\begin{proof}
Contravariance of dualization gives
\[
    (g\circ f)^\vee
    \simeq
    f^\vee\circ g^\vee,
\]
which proves the first identity. The identity morphism is immediate.

For involutivity, the Hermitian condition implies
\[
    h_M^\vee
    \simeq
    h_M
    \circ
    \operatorname{bid}_M^{-1},
\]
and therefore
\[
    (h_M^{-1})^\vee
    \simeq
    \operatorname{bid}_M
    \circ
    h_M^{-1}.
\]
Thus
\[
\begin{aligned}
    (f^\dagger)^\dagger
    &=
    h_N^{-1}
    \circ
    (f^\dagger)^\vee
    \circ
    h_M
    \\
    &\simeq
    h_N^{-1}
    \circ
    h_N^\vee
    \circ
    f^{\vee\vee}
    \circ
    (h_M^{-1})^\vee
    \circ
    h_M
    \\
    &\simeq
    \operatorname{bid}_N^{-1}
    \circ
    f^{\vee\vee}
    \circ
    \operatorname{bid}_M.
\end{aligned}
\]
Naturality of biduality gives
\[
    f^{\vee\vee}
    \circ
    \operatorname{bid}_M
    \simeq
    \operatorname{bid}_N
    \circ
    f.
\]
Hence
\[
    (f^\dagger)^\dagger
    \simeq
    f.
\]
\end{proof}

\section{Quadratic functors, cross-effects, and polarization}
\label{sec:quadratic-cross-effects-polarization}

\begin{definition}[Cross-effect and polarization]
\label{def:cross-effect-and-polarization}
Let
\[
    \mathcal C
\]
be a stable $\infty$-category, and let
\[
    \mathcal Q
    \colon
    \mathcal C^{\op}
    \longrightarrow
    \Sp
\]
be a reduced functor.

Its second cross-effect is
\[
    \mathbb B_{\mathcal Q}(X,Y)
    :=
    \fib
    \left(
        \mathcal Q(X\oplus Y)
        \longrightarrow
        \mathcal Q(X)\oplus\mathcal Q(Y)
    \right),
\]
where the two components are induced contravariantly by the canonical
inclusions
\[
    i_1
    \colon
    X
    \longrightarrow
    X\oplus Y,
\]
and
\[
    i_2
    \colon
    Y
    \longrightarrow
    X\oplus Y.
\]

Equivalently,
\[
    \mathbb B_{\mathcal Q}(X,Y)
\]
is characterized by the pullback square
\[
\begin{tikzcd}[row sep=large, column sep=large]
    \mathbb B_{\mathcal Q}(X,Y)
        \arrow[r]
        \arrow[d]
        \arrow[dr, phantom, very near start, "\lrcorner"]
    &
    \mathcal Q(X\oplus Y)
        \arrow[d, "{(\mathcal Q(i_1),\mathcal Q(i_2))}"]
    \\
    0
        \arrow[r]
    &
    \mathcal Q(X)\oplus\mathcal Q(Y).
\end{tikzcd}
\]

The symmetry
\[
    X\oplus Y
    \xrightarrow{\;\simeq\;}
    Y\oplus X
\]
canonically refines the cross-effect to a symmetric bifunctor. Thus there
is a coherent natural equivalence
\[
    \mathbb B_{\mathcal Q}(X,Y)
    \xrightarrow{\;\simeq\;}
    \mathbb B_{\mathcal Q}(Y,X)
\]
whose iterated transposition is coherently equivalent to the identity.

Consequently, the diagonal functor
\[
    X
    \longmapsto
    \mathbb B_{\mathcal Q}(X,X)
\]
is canonically a
\[
    C_2
\]
-object of
\[
    \Fun(\mathcal C^{\op},\Sp).
\]

The equivariant cross-effect construction supplies canonical natural
transformations
\[
    \mathbb B_{\mathcal Q}(X,X)_{hC_2}
    \longrightarrow
    \mathcal Q(X)
    \xrightarrow{\;\pol_X^{hC_2}\;}
    \mathbb B_{\mathcal Q}(X,X)^{hC_2}.
\]
See
\cite[Lemmas~1.1.9--1.1.10]{CalmesDottoHarpazEtAl}.

The second transformation is called the
\textbf{homotopy-fixed-point-valued polarization map}.

Its composite with the forgetful morphism
\[
    \mathbb B_{\mathcal Q}(X,X)^{hC_2}
    \longrightarrow
    \mathbb B_{\mathcal Q}(X,X)
\]
is the underlying polarization
\[
    \pol_X
    \colon
    \mathcal Q(X)
    \longrightarrow
    \mathbb B_{\mathcal Q}(X,X).
\]

The underlying polarization admits the following explicit description.
Let
\[
    \nabla
    \colon
    X\oplus X
    \longrightarrow
    X
\]
be the fold map, and let
\[
    p_1,p_2
    \colon
    X\oplus X
    \longrightarrow
    X
\]
be the two projections.

The additive structure on the mapping spectrum
\[
    \map_{\Sp}
    \bigl(
        \mathcal Q(X),
        \mathcal Q(X\oplus X)
    \bigr)
\]
defines a morphism
\[
\begin{aligned}
    a_X
    &:=
    \mathcal Q(\nabla)
    -
    \mathcal Q(p_1)
    -
    \mathcal Q(p_2)
    \\
    &\colon
    \mathcal Q(X)
    \longrightarrow
    \mathcal Q(X\oplus X).
\end{aligned}
\]

We verify that the composite of $a_X$ with each of the two maps defining
the cross-effect is canonically nullhomotopic.

Contravariance gives
\[
\begin{aligned}
    \mathcal Q(i_1)\circ a_X
    &=
    \mathcal Q(\nabla\circ i_1)
    -
    \mathcal Q(p_1\circ i_1)
    -
    \mathcal Q(p_2\circ i_1)
    \\
    &=
    \id_{\mathcal Q(X)}
    -
    \id_{\mathcal Q(X)}
    -
    \mathcal Q(0_{X,X}),
\end{aligned}
\]
where
\[
    0_{X,X}
    \colon
    X
    \longrightarrow
    X
\]
is the zero endomorphism.

The zero endomorphism factors through the zero object:
\[
    X
    \longrightarrow
    0
    \longrightarrow
    X.
\]
By contravariance,
\[
    \mathcal Q(0_{X,X})
    \colon
    \mathcal Q(X)
    \longrightarrow
    \mathcal Q(X)
\]
therefore factors through
\[
    \mathcal Q(0).
\]
Since $\mathcal Q$ is reduced,
\[
    \mathcal Q(0)
    \simeq
    0,
\]
so
\[
    \mathcal Q(0_{X,X})
\]
is canonically nullhomotopic. The first two terms cancel through the
additive structure on the mapping spectrum. Hence
\[
    \mathcal Q(i_1)\circ a_X
    \simeq
    0.
\]

The same argument gives
\[
    \mathcal Q(i_2)\circ a_X
    \simeq
    0.
\]

Together, these nullhomotopies determine a coherent nullhomotopy of
\[
\begin{aligned}
    \mathcal Q(X)
    &\xrightarrow{a_X}
    \mathcal Q(X\oplus X)
    \\
    &\xrightarrow{
        (\mathcal Q(i_1),\mathcal Q(i_2))
    }
    \mathcal Q(X)\oplus\mathcal Q(X).
\end{aligned}
\]

The universal property of the defining pullback square for
\[
    \mathbb B_{\mathcal Q}(X,X)
\]
therefore gives a morphism
\[
    \pol_X
    \colon
    \mathcal Q(X)
    \longrightarrow
    \mathbb B_{\mathcal Q}(X,X)
\]
whose composite with
\[
    \mathbb B_{\mathcal Q}(X,X)
    \longrightarrow
    \mathcal Q(X\oplus X)
\]
is
\[
    a_X.
\]

By
\cite[Lemma~1.1.10]{CalmesDottoHarpazEtAl},
this underlying polarization canonically refines, together with all
coherent equivariance data, to the homotopy-fixed-point-valued
polarization
\[
    \pol_X^{hC_2}
    \colon
    \mathcal Q(X)
    \longrightarrow
    \mathbb B_{\mathcal Q}(X,X)^{hC_2}.
\]

The preceding fold-map calculation identifies the composite
\[
    \mathcal Q(X)
    \xrightarrow{\pol_X^{hC_2}}
    \mathbb B_{\mathcal Q}(X,X)^{hC_2}
    \longrightarrow
    \mathbb B_{\mathcal Q}(X,X)
\]
with the factorization of
\[
    \mathcal Q(\nabla)
    -
    \mathcal Q(p_1)
    -
    \mathcal Q(p_2)
\]
through the second cross-effect.
\end{definition}

\section{Diagonals of biexact functors}
\label{sec:diagonal-biexact-functors}

\begin{lemma}[Diagonal of a biexact functor]
\label{lem:diagonal-biexact-two-excisive}
Let
\[
    \mathcal C
    \qquad\text{and}\qquad
    \mathcal D
\]
be stable $\infty$-categories, and let
\[
    B
    \colon
    \mathcal C\times\mathcal C
    \longrightarrow
    \mathcal D
\]
be exact separately in each variable. Then the diagonal functor
\[
    G(X)
    :=
    B(X,X)
\]
is reduced and
\[
    2
\]
-excisive.
\end{lemma}

\begin{proof}
Exactness in either variable implies
\[
    B(0,X)
    \simeq
    0
    \simeq
    B(X,0).
\]
Hence
\[
    G(0)
    =
    B(0,0)
    \simeq
    0,
\]
so
\[
    G
\]
is reduced.

Since
\[
    B
\]
is exact separately in each variable, it is
\[
    1
\]
-excisive separately in each variable. The multivariable diagonal theorem
for Goodwillie excision implies that the diagonal of a
\[
    (d_1,d_2)
\]
-excisive two-variable functor is
\[
    (d_1+d_2)
\]
-excisive. Applying it with
\[
    d_1=d_2=1
\]
shows that
\[
    G(X)=B(X,X)
\]
is
\[
    2
\]
-excisive.
\end{proof}


\section{Represented cross-effects and Poincar\'e structures}
\label{sec:represented-cross-effects-poincare}

\subsection{Canonical biduality from a represented cross-effect}

\begin{proposition}[Canonical biduality from a represented cross-effect]
\label{prop:canonical-biduality-from-cross-effect}
Let
\[
    \mathcal Q
    \colon
    \mathcal C^{\op}
    \longrightarrow
    \Sp
\]
be reduced and
\[
    2
\]
-excisive. Suppose that its symmetric cross-effect is right represented by
a functor
\[
    D_{\mathcal Q}
    \colon
    \mathcal C^{\op}
    \longrightarrow
    \mathcal C,
\]
in the sense that there are natural equivalences
\[
    \mathbb B_{\mathcal Q}(X,Y)
    \simeq
    \map_{\mathcal C}
    \bigl(
        X,
        D_{\mathcal Q}Y
    \bigr).
\]

Then
\[
    D_{\mathcal Q}
\]
is exact, and the involutive symmetry of the cross-effect determines an
adjunction
\[
    D_{\mathcal Q}^{\op}
    \colon
    \mathcal C
    \rightleftarrows
    \mathcal C^{\op}
    \colon
    D_{\mathcal Q}.
\]
Its unit is a canonical natural transformation
\[
    \operatorname{bid}
    \colon
    \id_{\mathcal C}
    \longrightarrow
    D_{\mathcal Q}D_{\mathcal Q}^{\op},
\]
or equivalently,
\[
    \operatorname{bid}_X
    \colon
    X
    \longrightarrow
    D_{\mathcal Q}^{2}X.
\]

Moreover, the following conditions are equivalent:
\begin{enumerate}[label=(\roman*)]
    \item every
    \[
        \operatorname{bid}_X
    \]
    is an equivalence;

    \item the adjunction
    \[
        D_{\mathcal Q}^{\op}
        \dashv
        D_{\mathcal Q}
    \]
    is an adjoint equivalence;

    \item
    \[
        D_{\mathcal Q}
        \colon
        \mathcal C^{\op}
        \longrightarrow
        \mathcal C
    \]
    is an equivalence.
\end{enumerate}
\end{proposition}

\begin{proof}
Since
\[
    \mathcal Q
\]
is reduced and
\[
    2
\]
-excisive, its cross-effect
\[
    \mathbb B_{\mathcal Q}
\]
is exact separately in each variable. For every
\[
    X\in\mathcal C,
\]
the functor
\[
    Y
    \longmapsto
    \map_{\mathcal C}
    \bigl(
        X,
        D_{\mathcal Q}Y
    \bigr)
\]
is therefore exact. The stable Yoneda embedding is exact and conservative,
so
\[
    D_{\mathcal Q}
\]
is exact.

Write
\[
    \theta_{X,Y}
    \colon
    \map_{\mathcal C}
    \bigl(
        X,
        D_{\mathcal Q}Y
    \bigr)
    \xrightarrow{\;\simeq\;}
    \map_{\mathcal C}
    \bigl(
        Y,
        D_{\mathcal Q}X
    \bigr)
\]
for the equivalence induced by the symmetry of the cross-effect. Its
coherence gives
\[
    \theta_{Y,X}
    \circ
    \theta_{X,Y}
    \simeq
    \id.
\]

For
\[
    X\in\mathcal C
\]
and
\[
    Y\in\mathcal C^{\op},
\]
one has
\[
\begin{aligned}
    \map_{\mathcal C^{\op}}
    \bigl(
        D_{\mathcal Q}^{\op}X,
        Y
    \bigr)
    &=
    \map_{\mathcal C}
    \bigl(
        Y,
        D_{\mathcal Q}X
    \bigr)
    \\
    &\simeq
    \map_{\mathcal C}
    \bigl(
        X,
        D_{\mathcal Q}Y
    \bigr).
\end{aligned}
\]
Thus
\[
    D_{\mathcal Q}^{\op}
    \dashv
    D_{\mathcal Q}.
\]

The unit
\[
    \operatorname{bid}_X
    \colon
    X
    \longrightarrow
    D_{\mathcal Q}^2X
\]
is characterized by
\[
    \theta_{X,D_{\mathcal Q}X}
    \bigl(
        \operatorname{bid}_X
    \bigr)
    \simeq
    \id_{D_{\mathcal Q}X}.
\]
Since
\[
    \theta_{X,D_{\mathcal Q}X}^{-1}
    \simeq
    \theta_{D_{\mathcal Q}X,X},
\]
one has
\[
    \operatorname{bid}_X
    =
    \theta_{D_{\mathcal Q}X,X}
    \bigl(
        \id_{D_{\mathcal Q}X}
    \bigr).
\]

Let
\[
    \epsilon_X
    \colon
    D_{\mathcal Q}^{\op}D_{\mathcal Q}X
    \longrightarrow
    X
\]
be the counit in
\[
    \mathcal C^{\op}.
\]
Its opposite is a morphism
\[
    \epsilon_X^{\op}
    \colon
    X
    \longrightarrow
    D_{\mathcal Q}^2X
\]
in
\[
    \mathcal C.
\]
The counit is characterized by
\[
    \theta_{X,D_{\mathcal Q}X}
    \bigl(
        \epsilon_X^{\op}
    \bigr)
    \simeq
    \id_{D_{\mathcal Q}X}.
\]
Hence
\[
    \epsilon_X^{\op}
    =
    \theta_{D_{\mathcal Q}X,X}
    \bigl(
        \id_{D_{\mathcal Q}X}
    \bigr)
    =
    \operatorname{bid}_X.
\]

Therefore the unit is an equivalence at every object if and only if the
counit is an equivalence at every object. This is equivalent to the
adjunction being an adjoint equivalence, and therefore to
\[
    D_{\mathcal Q}
    \colon
    \mathcal C^{\op}
    \longrightarrow
    \mathcal C
\]
being an equivalence.
\end{proof}


\subsection{Poincar\'e structures and Poincar\'e objects}

\begin{definition}[Poincar\'e structure]
\label{def:quadratic-poincare-structure}
Let
\[
    \mathcal C
\]
be a stable $\infty$-category. A \emph{Poincar\'e structure} on
\[
    \mathcal C
\]
is a reduced
\[
    2
\]
-excisive functor
\[
    \mathcal Q
    \colon
    \mathcal C^{\op}
    \longrightarrow
    \Sp
\]
such that:
\begin{enumerate}[label=(\roman*)]
    \item the cross-effect
    \[
        \mathbb B_{\mathcal Q}(X,Y)
        =
        \fib
        \left(
            \mathcal Q(X\oplus Y)
            \longrightarrow
            \mathcal Q(X)\oplus\mathcal Q(Y)
        \right)
    \]
    is right represented by an exact functor
    \[
        D_{\mathcal Q}
        \colon
        \mathcal C^{\op}
        \longrightarrow
        \mathcal C;
    \]

    \item the canonical biduality transformation of
    \Ref{prop:canonical-biduality-from-cross-effect},
    \[
        \operatorname{bid}_X
        \colon
        X
        \longrightarrow
        D_{\mathcal Q}^{2}X,
    \]
    is an equivalence for every
    \[
        X\in\mathcal C.
    \]
\end{enumerate}

For
\[
    q\in\Omega^\infty\mathcal Q(X),
\]
define
\[
    q^\sharp
    \colon
    X
    \longrightarrow
    D_{\mathcal Q}X
\]
to be the morphism corresponding, under
\[
    \Omega^\infty
    \mathbb B_{\mathcal Q}(X,X)
    \simeq
    \Map_{\mathcal C}
    \bigl(
        X,
        D_{\mathcal Q}X
    \bigr),
\]
to the underlying polarized class
\[
    \pol_X(q)
    \in
    \Omega^\infty
    \mathbb B_{\mathcal Q}(X,X).
\]

A \emph{Poincar\'e object} is a pair
\[
    (X,q)
\]
with
\[
    q\in\Omega^\infty\mathcal Q(X)
\]
such that
\[
    q^\sharp
    \colon
    X
    \xrightarrow{\;\simeq\;}
    D_{\mathcal Q}X
\]
is an equivalence.
\end{definition}


\subsection{Spaces of forms and Poincar\'e objects}

\begin{definition}[Spaces of forms and Poincar\'e objects]
\label{def:spaces-of-poincare-objects}
Let
\[
    \mathcal Q
    \colon
    \mathcal C^{\op}
    \longrightarrow
    \Sp
\]
be a Poincar\'e structure.

Let
\[
    \operatorname{He}(\mathcal C,\mathcal Q)
    \longrightarrow
    \mathcal C
\]
be the right fibration classified by
\[
    \Omega^\infty\mathcal Q
    \colon
    \mathcal C^{\op}
    \longrightarrow
    \mathcal S.
\]
Thus an object is a pair
\[
    (X,q),
    \qquad
    q\in\Omega^\infty\mathcal Q(X).
\]

For a morphism
\[
    f
    \colon
    X
    \longrightarrow
    Y,
\]
a morphism from
\[
    (X,q_X)
\]
to
\[
    (Y,q_Y)
\]
lying over
\[
    f
\]
is represented by a path
\[
    q_X
    \simeq
    \mathcal Q(f)(q_Y)
\]
in
\[
    \Omega^\infty\mathcal Q(X),
\]
together with the higher coherent data encoded by the right fibration.

Define the \emph{space of forms} by
\[
    \operatorname{Fm}(\mathcal C,\mathcal Q)
    :=
    \operatorname{He}(\mathcal C,\mathcal Q)^\simeq.
\]

The \emph{space of Poincar\'e objects}
\[
    \operatorname{Pn}(\mathcal C,\mathcal Q)
    \subseteq
    \operatorname{Fm}(\mathcal C,\mathcal Q)
\]
is the full $\infty$-subgroupoid spanned by the pairs
\[
    (X,q)
\]
for which
\[
    q^\sharp
    \colon
    X
    \xrightarrow{\;\simeq\;}
    D_{\mathcal Q}X
\]
is an equivalence.
\end{definition}


\subsection{Perfect cross-effects and quadratic terminology}

\begin{remark}[Comparison with perfect bilinear cross-effects]
\label{rem:poincare-perfect-bilinear-cross-effect}
The preceding definition is equivalent to the usual requirement that the
symmetric bilinear cross-effect
\[
    \mathbb B_{\mathcal Q}
\]
be perfect.

Indeed, right representability produces the self-adjunction
\[
    D_{\mathcal Q}^{\op}
    \dashv
    D_{\mathcal Q},
\]
and invertibility of the canonical biduality transformation makes this
adjunction an adjoint equivalence. Thus
\[
    D_{\mathcal Q}
    \colon
    \mathcal C^{\op}
    \xrightarrow{\;\simeq\;}
    \mathcal C
\]
is the perfect duality represented by the cross-effect.
\end{remark}

\begin{remark}[Goodwillie-quadratic versus homotopy-quadratic]
\label{rem:goodwillie-quadratic-terminology}
The condition that
\[
    \mathcal Q
\]
is reduced and
\[
    2
\]
-excisive is frequently described by saying that
\[
    \mathcal Q
\]
is quadratic in the Goodwillie sense. This should not be confused with the
specific homotopy-quadratic refinement
\[
    \mathcal Q^{\mathrm q}(X)
    =
    \mathbb B(X,X)_{hC_2}.
\]
\end{remark}


\section{The homotopy-symmetric refinement of a coherent duality}
\label{sec:homotopy-symmetric-refinement-general}

Let
\[
    \mathcal C
\]
be a stable $\infty$-category equipped with an exact coherent perfect
duality. The following proposition constructs the canonical
homotopy-symmetric refinement and identifies its Poincar\'e objects before
any specialization to module or quasi-coherent categories.

\begin{proposition}[Hermitian objects as Poincar\'e objects]
\label{prop:hermitian-as-poincare}
Let
\[
    \mathcal C
\]
be a stable $\infty$-category equipped with an exact coherent perfect
duality
\[
    D
    \colon
    \mathcal C^{\op}
    \xrightarrow{\;\simeq\;}
    \mathcal C
\]
and biduality
\[
    \operatorname{bid}_X
    \colon
    X
    \xrightarrow{\;\simeq\;}
    D^2X.
\]
Define
\[
    \mathbb B(X,Y)
    :=
    \map_{\mathcal C}(X,DY),
\]
with transposition
\[
    \mathbb B(X,Y)
    \xrightarrow{\;\simeq\;}
    \mathbb B(Y,X)
\]
sending
\[
    f
    \colon
    X
    \longrightarrow
    DY
\]
to
\[
    Y
    \xrightarrow{\operatorname{bid}_Y}
    D^2Y
    \xrightarrow{D(f)}
    DX.
\]
Define
\[
    \mathcal Q^{\mathrm s}(X)
    :=
    \mathbb B(X,X)^{hC_2}.
\]

Then:
\begin{enumerate}[label=(\roman*)]
    \item
    \[
        \mathcal Q^{\mathrm s}
    \]
    is reduced and
    \[
        2
    \]
    -excisive;

    \item its cross-effect is canonically equivalent to
    \[
        \mathbb B:
    \]
    \[
        \mathbb B_{\mathcal Q^{\mathrm s}}(X,Y)
        \simeq
        \mathbb B(X,Y);
    \]

    \item under the equivalence
    \[
        \mathbb B_{\mathcal Q^{\mathrm s}}(X,X)
        \simeq
        \mathbb B(X,X),
    \]
    the homotopy-fixed-point-valued polarization map identifies with the
    identity
    \[
        \mathbb B(X,X)^{hC_2}
        \longrightarrow
        \mathbb B(X,X)^{hC_2}.
    \]
    Consequently, its underlying polarization identifies with the
    forgetful map
    \[
        \mathbb B(X,X)^{hC_2}
        \longrightarrow
        \mathbb B(X,X);
    \]

    \item there is a canonical equivalence
    \[
        \operatorname{Pn}
        \bigl(
            \mathcal C,
            \mathcal Q^{\mathrm s}
        \bigr)
        \xrightarrow{\;\simeq\;}
        \bigl(
            \mathcal C^\simeq
        \bigr)^{hC_2}.
    \]
\end{enumerate}

Equivalently, the
\[
    \mathcal Q^{\mathrm s}
\]
-Poincar\'e objects are precisely the coherently symmetric equivalences
\[
    h
    \colon
    X
    \xrightarrow{\;\simeq\;}
    DX.
\]
\end{proposition}
\begin{proof}
Because
\[
    D
    \colon
    \mathcal C^{\op}
    \longrightarrow
    \mathcal C
\]
is exact, the functor
\[
    \mathbb B
    \colon
    \mathcal C^{\op}\times\mathcal C^{\op}
    \longrightarrow
    \Sp
\]
is exact separately in each variable.

Apply
\Ref{lem:diagonal-biexact-two-excisive}
to the stable $\infty$-category
\[
    \mathcal C^{\op}
\]
and the biexact functor
\[
    \mathbb B
    \colon
    \mathcal C^{\op}\times\mathcal C^{\op}
    \longrightarrow
    \Sp.
\]
It follows that the diagonal
\[
    X
    \longmapsto
    \mathbb B(X,X)
\]
is reduced and $2$-excisive.

The coherent transposition equips this diagonal with a
\[
    C_2
\]
-action. Limits in
\[
    \Sp^{BC_2}
\]
are created in
\[
    \Sp.
\]
Consequently the equivariant diagonal is again $2$-excisive as an
\[
    \Sp^{BC_2}
\]
-valued functor.

Homotopy fixed points are computed by the limit functor
\[
    (-)^{hC_2}
    \colon
    \Sp^{BC_2}
    \longrightarrow
    \Sp.
\]
Since limits preserve cartesian cubes, it follows that
\[
    \mathcal Q^{\mathrm s}(X)
    =
    \mathbb B(X,X)^{hC_2}
\]
is reduced and $2$-excisive.

We next compute its cross-effect. Bilinearity gives a natural
decomposition
\[
\begin{aligned}
    \mathbb B(X\oplus Y,X\oplus Y)
    \simeq{}&
    \mathbb B(X,X)
    \oplus
    \mathbb B(X,Y)
    \\
    &\oplus
    \mathbb B(Y,X)
    \oplus
    \mathbb B(Y,Y).
\end{aligned}
\]
The $C_2$-action preserves the two diagonal summands and exchanges the two
off-diagonal summands through transposition. Since homotopy fixed points
preserve fibres,
\[
\begin{aligned}
    \mathbb B_{\mathcal Q^{\mathrm s}}(X,Y)
    &\simeq
    \left(
        \mathbb B(X,Y)
        \oplus
        \mathbb B(Y,X)
    \right)^{hC_2}.
\end{aligned}
\]

Using the transposition equivalence
\[
    \mathbb B(X,Y)
    \xrightarrow{\;\simeq\;}
    \mathbb B(Y,X),
\]
the off-diagonal $C_2$-spectrum is canonically equivalent to
\[
    \operatorname{Coind}_{1}^{C_2}
    \mathbb B(X,Y).
\]
Homotopy fixed points of a coinduced spectrum recover the original
spectrum. Hence
\[
    \left(
        \mathbb B(X,Y)
        \oplus
        \mathbb B(Y,X)
    \right)^{hC_2}
    \simeq
    \mathbb B(X,Y).
\]
This equivalence is natural in $X$ and $Y$ and is compatible with
transposition. Thus there is an equivalence of symmetric bilinear
functors
\[
    \mathbb B_{\mathcal Q^{\mathrm s}}
    \xrightarrow{\;\simeq\;}
    \mathbb B.
\]

We now identify the homotopy-fixed-point-valued polarization.

By
\cite[Corollary~1.3.6]{CalmesDottoHarpazEtAl},
the symmetric cross-effect functor
\[
    \mathbb B_{(-)}
    \colon
    \Fun^{\mathrm q}(\mathcal C)
    \longrightarrow
    \Fun^{\mathrm s}(\mathcal C)
\]
admits a fully faithful right adjoint
\[
    \mathbb B
    \longmapsto
    \mathcal Q_{\mathbb B}^{\mathrm s},
    \qquad
    \mathcal Q_{\mathbb B}^{\mathrm s}(X)
    :=
    \mathbb B(X,X)^{hC_2}.
\]
The unit of this adjunction at a reduced quadratic functor $\mathcal Q$ is
the homotopy-fixed-point-valued polarization
\[
    \mathcal Q(X)
    \longrightarrow
    \mathbb B_{\mathcal Q}(X,X)^{hC_2}.
\]

For
\[
    \mathcal Q^{\mathrm s}
    =
    \mathcal Q_{\mathbb B}^{\mathrm s},
\]
the counit of the adjunction is the cross-effect equivalence
\[
    \mathbb B_{\mathcal Q^{\mathrm s}}
    \xrightarrow{\;\simeq\;}
    \mathbb B
\]
constructed above. Under this equivalence, the triangle identity for the
adjunction identifies the unit at
\[
    \mathcal Q_{\mathbb B}^{\mathrm s}
\]
with the identity transformation
\[
    \mathbb B(X,X)^{hC_2}
    \longrightarrow
    \mathbb B(X,X)^{hC_2}.
\]
Therefore the homotopy-fixed-point-valued polarization of
\[
    \mathcal Q^{\mathrm s}
\]
identifies with the identity.

After forgetting the homotopy-fixed-point structure, the underlying
polarization identifies with the forgetful morphism
\[
    \mathbb B(X,X)^{hC_2}
    \longrightarrow
    \mathbb B(X,X).
\]

It remains to identify the Poincaré objects.

Since
\[
    \Omega^\infty
    \colon
    \Sp
    \longrightarrow
    \mathcal S
\]
is a right adjoint, it preserves limits. Hence there is a canonical
equivalence
\[
    \Omega^\infty
    \left(
        \mathbb B(X,X)^{hC_2}
    \right)
    \simeq
    \left(
        \Omega^\infty\mathbb B(X,X)
    \right)^{hC_2}.
\]

Using
\[
    \mathbb B(X,X)
    =
    \map_{\mathcal C}(X,DX),
\]
a point
\[
    q
    \in
    \Omega^\infty\mathcal Q^{\mathrm s}(X)
\]
therefore consists of a morphism
\[
    q^\sharp
    \colon
    X
    \longrightarrow
    DX
\]
together with a coherent homotopy-fixed-point structure for
transposition. Its first fixed-point coherence is
\[
    D(q^\sharp)
    \circ
    \operatorname{bid}_X
    \simeq
    q^\sharp.
\]

The pair
\[
    (X,q)
\]
is a Poincaré object precisely when
\[
    q^\sharp
    \colon
    X
    \xrightarrow{\;\simeq\;}
    DX
\]
is an equivalence.

The coherent duality determines a covariant $C_2$-action on
\[
    \mathcal C^\simeq
\]
which sends an object $X$ to $DX$ and sends an equivalence
\[
    f
    \colon
    X
    \xrightarrow{\;\simeq\;}
    Y
\]
to
\[
    D(f^{-1})
    \colon
    DX
    \xrightarrow{\;\simeq\;}
    DY.
\]
A homotopy-fixed-point object for this action is therefore an object $X$
equipped with an equivalence
\[
    h
    \colon
    X
    \xrightarrow{\;\simeq\;}
    DX
\]
and coherent fixed-point data. Its first cocycle condition is
\[
    D(h)
    \circ
    \operatorname{bid}_X
    \simeq
    h.
\]

The correspondence also agrees on equivalences. Let
\[
    (X,q_X)
    \qquad\text{and}\qquad
    (Y,q_Y)
\]
be Poincaré objects. An equivalence
\[
    f
    \colon
    X
    \xrightarrow{\;\simeq\;}
    Y
\]
preserves the Poincaré forms precisely when
\[
    q_X
    \simeq
    \mathcal Q^{\mathrm s}(f)(q_Y).
\]
After forgetting the homotopy-fixed-point structures, this condition is
\[
    q_X^\sharp
    \simeq
    D(f)
    \circ
    q_Y^\sharp
    \circ
    f.
\]

On the other hand, the condition that $f$ be an equivalence between the
corresponding homotopy-fixed-point objects is
\[
    q_Y^\sharp
    \circ
    f
    \simeq
    D(f^{-1})
    \circ
    q_X^\sharp.
\]
Applying $D(f)$ identifies this condition with
\[
    q_X^\sharp
    \simeq
    D(f)
    \circ
    q_Y^\sharp
    \circ
    f.
\]

Finally,
\cite[Proposition~2.2.11]{CalmesDottoHarpazEtAl}
identifies the full coherent homotopy types and supplies a canonical
equivalence of $\infty$-groupoids
\[
    \operatorname{Pn}
    \bigl(
        \mathcal C,
        \mathcal Q^{\mathrm s}
    \bigr)
    \xrightarrow{\;\simeq\;}
    \bigl(
        \mathcal C^\simeq
    \bigr)^{hC_2}.
\]
\end{proof}


\section{The canonical homotopy-symmetric Poincar\'e structure on perfect objects}
\label{sec:canonical-homotopy-symmetric-poincare}

\subsection{The affine bilinear pairing and its transposition}

Retain the notation
\[
    \mathcal C_A
    :=
    \Perf(A),
    \qquad
    D_A
    :=
    (-)^\vee.
\]

Under
\Ref{prop:canonical-bilinear-symmetry-representability},
the canonical bilinear functor may equally be written
\[
    \mathbb B_A(M,N)
    \simeq
    \map_{\Perf(A)}
    \bigl(
        M,
        N^\vee
    \bigr).
\]

\begin{proposition}[Symmetry of the bilinear mapping spectrum]
\label{prop:bilinear-mapping-spectrum-symmetry}
There is a natural coherent symmetry equivalence
\[
    \sigma_{M,N}
    \colon
    \mathbb B_A(M,N)
    \xrightarrow{\;\simeq\;}
    \mathbb B_A(N,M)
\]
which, under the represented description, sends a morphism
\[
    q
    \colon
    M
    \longrightarrow
    N^\vee
\]
to
\[
    \sigma_{M,N}(q)
    :=
    q^\vee
    \circ
    \operatorname{bid}_N
    \colon
    N
    \longrightarrow
    M^\vee.
\]
The composite
\[
    \sigma_{N,M}
    \circ
    \sigma_{M,N}
\]
is canonically equivalent to the identity, together with all higher
\[
    C_2
\]
-coherences.
\end{proposition}

\begin{proof}
The formula was established in the proof of
\Ref{prop:canonical-perfect-duality}.
For
\[
    q
    \colon
    M
    \longrightarrow
    N^\vee,
\]
one has
\[
\begin{aligned}
    \sigma_{N,M}
    \sigma_{M,N}(q)
    &=
    \bigl(
        q^\vee
        \operatorname{bid}_N
    \bigr)^\vee
    \operatorname{bid}_M
    \\
    &\simeq
    \operatorname{bid}_N^\vee
    \circ
    q^{\vee\vee}
    \circ
    \operatorname{bid}_M
    \\
    &\simeq
    \operatorname{bid}_N^\vee
    \circ
    \operatorname{bid}_{N^\vee}
    \circ
    q
    \\
    &\simeq
    q.
\end{aligned}
\]
The third line is naturality of biduality, and the final line is the first
nontrivial coherence of the coherent perfect duality.

The higher coherences are inherited from the coherent symmetry of
\[
    \mathbb B_A.
\]
\end{proof}

On the diagonal, the symmetry defines a coherent
\[
    C_2
\]
-action on
\[
    \mathbb B_A(M,M)
    =
    \map_{\Perf(A)}(M,M^\vee).
\]
The nontrivial element acts by transposition:
\[
    \tau_M(q)
    :=
    q^\vee
    \circ
    \operatorname{bid}_M.
\]


\subsection{The affine homotopy-symmetric functor and its cross-effect}

\begin{definition}[Homotopy-symmetric Poincar\'e functor]
\label{def:homotopy-symmetric-poincare-functor}
Define
\[
    \mathcal Q_A^{\mathrm s}(M)
    :=
    \mathbb B_A(M,M)^{hC_2}.
\]

For a morphism
\[
    f
    \colon
    M
    \longrightarrow
    N,
\]
define
\[
    f^*
    \colon
    \mathbb B_A(N,N)
    \longrightarrow
    \mathbb B_A(M,M)
\]
by
\[
    f^*(q)
    :=
    f^\vee
    \circ
    q
    \circ
    f.
\]
The induced map on homotopy fixed points is denoted
\[
    \mathcal Q_A^{\mathrm s}(f)
    \colon
    \mathcal Q_A^{\mathrm s}(N)
    \longrightarrow
    \mathcal Q_A^{\mathrm s}(M).
\]
\end{definition}

\begin{lemma}[Equivariance of pullback]
\label{lem:poincare-pullback-equivariant}
For every morphism
\[
    f
    \colon
    M
    \longrightarrow
    N,
\]
the map
\[
    f^*
    \colon
    \mathbb B_A(N,N)
    \longrightarrow
    \mathbb B_A(M,M)
\]
is
\[
    C_2
\]
-equivariant.
\end{lemma}

\begin{proof}
For
\[
    q
    \colon
    N
    \longrightarrow
    N^\vee,
\]
one has
\[
\begin{aligned}
    \tau_M(f^*q)
    &=
    \bigl(
        f^\vee
        q
        f
    \bigr)^\vee
    \operatorname{bid}_M
    \\
    &\simeq
    f^\vee
    q^\vee
    f^{\vee\vee}
    \operatorname{bid}_M
    \\
    &\simeq
    f^\vee
    q^\vee
    \operatorname{bid}_N
    f
    \\
    &=
    f^*(\tau_Nq).
\end{aligned}
\]
The third line is naturality of biduality. The equality is compatible with
the full coherent
\[
    C_2
\]
-actions because all structure maps arise from the coherent perfect
duality.
\end{proof}

Consequently,
\[
    \mathcal Q_A^{\mathrm s}
    \colon
    \Perf(A)^{\op}
    \longrightarrow
    \Sp
\]
is a well-defined functor.

\begin{proposition}[Cross-effect of the homotopy-symmetric functor]
\label{prop:homotopy-symmetric-cross-effect}
The functor
\[
    \mathcal Q_A^{\mathrm s}
\]
is reduced and
\[
    2
\]
-excisive. Its symmetric bilinear cross-effect is naturally equivalent to
\[
    \mathbb B_A:
\]
\[
    \operatorname{cr}_2
    \mathcal Q_A^{\mathrm s}(M,N)
    \simeq
    \mathbb B_A(M,N).
\]
\end{proposition}

\begin{proof}
The functor is reduced because
\[
    \mathbb B_A(0,0)
    \simeq
    0.
\]

Apply
\Ref{lem:diagonal-biexact-two-excisive}
with
\[
    \mathcal C
    =
    \Perf(A)^{\op},
    \qquad
    \mathcal D
    =
    \Sp,
    \qquad
    B
    =
    \mathbb B_A.
\]
The underlying spectrum-valued diagonal
\[
    M
    \longmapsto
    \mathbb B_A(M,M)
\]
is
\[
    2
\]
-excisive. Since limits in
\[
    \Sp^{BC_2}
\]
are created in
\[
    \Sp,
\]
its coherent
\[
    C_2
\]
-equivariant refinement is likewise
\[
    2
\]
-excisive as an
\[
    \Sp^{BC_2}
\]
-valued functor.

Homotopy fixed points are computed by a limit
\[
    (-)^{hC_2}
    \colon
    \Sp^{BC_2}
    \longrightarrow
    \Sp.
\]
Limits preserve cartesian cubes. Hence
\[
    \mathcal Q_A^{\mathrm s}(M)
    =
    \mathbb B_A(M,M)^{hC_2}
\]
is
\[
    2
\]
-excisive.

For the cross-effect, bilinearity gives a natural decomposition
\[
\begin{aligned}
    \mathbb B_A(M\oplus N,M\oplus N)
    \simeq{}&
    \mathbb B_A(M,M)
    \oplus
    \mathbb B_A(M,N)
    \\
    &\oplus
    \mathbb B_A(N,M)
    \oplus
    \mathbb B_A(N,N).
\end{aligned}
\]
The
\[
    C_2
\]
-action preserves the two diagonal terms and exchanges the two
off-diagonal terms through
\[
    \sigma_{M,N}
    \colon
    \mathbb B_A(M,N)
    \xrightarrow{\;\simeq\;}
    \mathbb B_A(N,M).
\]

The cross-effect is the fibre of
\[
    \mathcal Q_A^{\mathrm s}(M\oplus N)
    \longrightarrow
    \mathcal Q_A^{\mathrm s}(M)
    \oplus
    \mathcal Q_A^{\mathrm s}(N).
\]
Since homotopy fixed points preserve fibres,
\[
\begin{aligned}
    \operatorname{cr}_2
    \mathcal Q_A^{\mathrm s}(M,N)
    &\simeq
    \bigl(
        \mathbb B_A(M,N)
        \oplus
        \mathbb B_A(N,M)
    \bigr)^{hC_2}.
\end{aligned}
\]
Using
\[
    \sigma_{M,N},
\]
the off-diagonal
\[
    C_2
\]
-spectrum is canonically equivalent to
\[
    \operatorname{Coind}_{1}^{C_2}
    \mathbb B_A(M,N).
\]
The homotopy fixed points of a coinduced spectrum recover the original
spectrum. Therefore
\[
    \bigl(
        \mathbb B_A(M,N)
        \oplus
        \mathbb B_A(N,M)
    \bigr)^{hC_2}
    \simeq
    \mathbb B_A(M,N).
\]
The construction is compatible with the symmetry interchanging
\[
    M
    \qquad\text{and}\qquad
    N,
\]
so this is an equivalence of symmetric bilinear functors.
\end{proof}


\subsection{Perfect modules}

\begin{corollary}[The canonical Poincar\'e structure on perfect modules]
\label{cor:canonical-poincare-perfect-modules}
The functor
\[
    \mathcal Q_A^{\mathrm s}
    \colon
    \Perf(A)^{\op}
    \longrightarrow
    \Sp
\]
is a Poincar\'e structure on
\[
    \Perf(A).
\]
Its represented cross-effect is
\[
    \mathbb B_A(M,N)
    \simeq
    \map_{\Perf(A)}(M,D_A N),
\]
and the representing duality is the canonical coherent rigid duality
\[
    D_A
    =
    (-)^\vee.
\]

There is a canonical equivalence
\[
    \operatorname{Pn}
    \bigl(
        \Perf(A),
        \mathcal Q_A^{\mathrm s}
    \bigr)
    \simeq
    \bigl(
        \Perf(A)^\simeq
    \bigr)^{hC_2}.
\]
Hence its objects are precisely the Hermitian perfect modules of
\Ref{def:hermitian-perfect-module}.
\end{corollary}

\begin{proof}
By
\Ref{prop:homotopy-symmetric-cross-effect},
the functor
\[
    \mathcal Q_A^{\mathrm s}
\]
is reduced and
\[
    2
\]
-excisive, with symmetric cross-effect
\[
    \mathbb B_A.
\]
By
\Ref{prop:canonical-bilinear-symmetry-representability},
\[
    \mathbb B_A(M,N)
    \simeq
    \map_{\Perf(A)}(M,D_A N).
\]
By
\Ref{prop:canonical-perfect-duality},
\[
    D_A
\]
is an exact coherent contravariant equivalence. The conclusion follows
from
\Ref{prop:hermitian-as-poincare}.
\end{proof}


\subsection{Perfect objects over a presheaf}

\begin{proposition}[Canonical Poincar\'e structure on perfect objects over a presheaf]
\label{prop:canonical-poincare-perfect-presheaves}
Let
\[
    X\in\PSh(\Aff),
\]
and let
\[
    \mathbf 1_X
\]
denote the tensor unit of
\[
    \Perf(X).
\]
Define
\[
    \mathbb B_X
    \colon
    \Perf(X)^{\op}
    \times
    \Perf(X)^{\op}
    \longrightarrow
    \Sp
\]
by
\[
    \mathbb B_X(M,N)
    :=
    \map_{\Perf(X)}
    \bigl(
        M\otimes N,
        \mathbf 1_X
    \bigr).
\]

Then
\[
    \mathbb B_X
\]
is exact separately in each variable and carries a canonical coherent
symmetry
\[
    \mathbb B_X(M,N)
    \xrightarrow{\;\simeq\;}
    \mathbb B_X(N,M).
\]
It is right-represented by the canonical rigid duality
\[
    D_X
    :=
    (-)^\vee
    \colon
    \Perf(X)^{\op}
    \xrightarrow{\;\simeq\;}
    \Perf(X),
\]
through natural equivalences
\[
    \mathbb B_X(M,N)
    \simeq
    \map_{\Perf(X)}
    \bigl(
        M,
        N^\vee
    \bigr).
\]

Define
\[
    \mathcal Q_X^{\mathrm s}(M)
    :=
    \mathbb B_X(M,M)^{hC_2},
\]
where the
\[
    C_2
\]
-action is induced by the coherent symmetry of
\[
    \mathbb B_X.
\]
Then
\[
    \mathcal Q_X^{\mathrm s}
    \colon
    \Perf(X)^{\op}
    \longrightarrow
    \Sp
\]
is a Poincar\'e structure on
\[
    \Perf(X).
\]
Its symmetric bilinear cross-effect is naturally equivalent to
\[
    \mathbb B_X,
\]
and its representing duality is
\[
    D_X=(-)^\vee.
\]

There is a canonical equivalence
\[
    \operatorname{Pn}
    \bigl(
        \Perf(X),
        \mathcal Q_X^{\mathrm s}
    \bigr)
    \xrightarrow{\;\simeq\;}
    \bigl(
        \Perf(X)^\simeq
    \bigr)^{hC_2}.
\]
Thus the
\[
    \mathcal Q_X^{\mathrm s}
\]
-Poincar\'e objects are precisely the coherently symmetric equivalences
\[
    M
    \xrightarrow{\;\simeq\;}
    M^\vee.
\]

\end{proposition}

\begin{proof}
By
\Ref{prop:perfect-stable-monoidal},
the category
\[
    \Perf(X)
\]
is a stable idempotent-complete rigid symmetric monoidal
$\infty$-category whose tensor product is exact separately in each
variable.

Tensoring is exact separately in each variable, and mapping into the
tensor unit
\[
    \mathbf 1_X
\]
is exact contravariantly. Hence
\[
    \mathbb B_X
\]
is exact separately in each variable.

The braiding
\[
    N\otimes M
    \xrightarrow{\;\simeq\;}
    M\otimes N
\]
induces a natural equivalence
\[
    \mathbb B_X(M,N)
    \xrightarrow{\;\simeq\;}
    \mathbb B_X(N,M).
\]
The symmetric monoidal coherence of the braiding supplies this equivalence
with its coherent involutivity and all higher
\[
    C_2
\]
-compatibilities.

Since every object of
\[
    \Perf(X)
\]
is dualizable, rigidity gives natural equivalences
\[
\begin{aligned}
    \mathbb B_X(M,N)
    &=
    \map_{\Perf(X)}
    \bigl(
        M\otimes N,
        \mathbf 1_X
    \bigr)
    \\
    &\simeq
    \map_{\Perf(X)}
    \bigl(
        M,
        N^\vee
    \bigr).
\end{aligned}
\]
Thus
\[
    \mathbb B_X
\]
is right-represented by
\[
    D_X=(-)^\vee.
\]

The canonical biduality morphism
\[
    M
    \longrightarrow
    M^{\vee\vee}
\]
is an equivalence for every
\[
    M\in\Perf(X).
\]
Therefore
\[
    D_X
    \colon
    \Perf(X)^{\op}
    \xrightarrow{\;\simeq\;}
    \Perf(X)
\]
is an exact coherent perfect duality, and
\[
    \mathbb B_X
\]
is a perfect symmetric bilinear functor.

Applying
\Ref{prop:hermitian-as-poincare}
to
\[
    \mathcal C=\Perf(X),
    \qquad
    D=D_X,
\]
shows that
\[
    \mathcal Q_X^{\mathrm s}(M)
    =
    \mathbb B_X(M,M)^{hC_2}
\]
is reduced and
\[
    2
\]
-excisive, that its cross-effect is naturally equivalent to
\[
    \mathbb B_X,
\]
and that it defines a Poincar\'e structure represented by
\[
    D_X.
\]
The same proposition supplies the canonical equivalence
\[
    \operatorname{Pn}
    \bigl(
        \Perf(X),
        \mathcal Q_X^{\mathrm s}
    \bigr)
    \xrightarrow{\;\simeq\;}
    \bigl(
        \Perf(X)^\simeq
    \bigr)^{hC_2}.
\]
\end{proof}

\section{Pullback of Hermitian and Poincar\'e objects}
\label{sec:pullback-hermitian-poincare}

\begin{proposition}[Functoriality of homotopy-symmetric Poincaré structures]
\label{prop:functoriality-homotopy-symmetric-structures}
Let
\[
    (\mathcal C,D_{\mathcal C})
    \qquad\text{and}\qquad
    (\mathcal D,D_{\mathcal D})
\]
be stable $\infty$-categories equipped with exact coherent perfect
dualities. Let
\[
    F
    \colon
    \mathcal C
    \longrightarrow
    \mathcal D
\]
be an exact duality-preserving functor, with coherent natural equivalence
\[
    \chi_X
    \colon
    F(D_{\mathcal C}X)
    \xrightarrow{\;\simeq\;}
    D_{\mathcal D}F(X).
\]

For
\[
    X,Y\in\mathcal C,
\]
there is a natural morphism
\[
\begin{aligned}
    \mathbb B_{\mathcal C}(X,Y)
    &=
    \map_{\mathcal C}
    \bigl(
        X,
        D_{\mathcal C}Y
    \bigr)
    \\
    &\longrightarrow
    \map_{\mathcal D}
    \bigl(
        F(X),
        D_{\mathcal D}F(Y)
    \bigr)
    =
    \mathbb B_{\mathcal D}
    \bigl(
        F(X),
        F(Y)
    \bigr),
\end{aligned}
\]
given by
\[
    q
    \longmapsto
    \chi_Y\circ F(q).
\]
These morphisms define a morphism of symmetric bilinear functors and are
compatible with transposition.

On the diagonal they induce a natural transformation
\[
    \mathcal Q_{\mathcal C}^{\mathrm s}
    \longrightarrow
    \mathcal Q_{\mathcal D}^{\mathrm s}
    \circ
    F^{\op}.
\]

Consequently, $F$ induces functors
\[
    \operatorname{He}
    \bigl(
        \mathcal C,
        \mathcal Q_{\mathcal C}^{\mathrm s}
    \bigr)
    \longrightarrow
    \operatorname{He}
    \bigl(
        \mathcal D,
        \mathcal Q_{\mathcal D}^{\mathrm s}
    \bigr),
\]
\[
    \operatorname{Fm}
    \bigl(
        \mathcal C,
        \mathcal Q_{\mathcal C}^{\mathrm s}
    \bigr)
    \longrightarrow
    \operatorname{Fm}
    \bigl(
        \mathcal D,
        \mathcal Q_{\mathcal D}^{\mathrm s}
    \bigr),
\]
and
\[
    \operatorname{Pn}
    \bigl(
        \mathcal C,
        \mathcal Q_{\mathcal C}^{\mathrm s}
    \bigr)
    \longrightarrow
    \operatorname{Pn}
    \bigl(
        \mathcal D,
        \mathcal Q_{\mathcal D}^{\mathrm s}
    \bigr).
\]

For a Poincaré object
\[
    (X,q),
\]
the underlying adjoint of the transported form is
\[
    F(X)
    \xrightarrow{F(q^\sharp)}
    F(D_{\mathcal C}X)
    \xrightarrow{\chi_X}
    D_{\mathcal D}F(X).
\]

The resulting functors are coherently compatible with identities and
composition of duality-preserving functors.
\end{proposition}

\begin{proof}
Applying $F$ to morphisms gives natural maps
\[
\begin{aligned}
    \map_{\mathcal C}
    \bigl(
        X,
        D_{\mathcal C}Y
    \bigr)
    \longrightarrow
    \map_{\mathcal D}
    \bigl(
        F(X),
        F(D_{\mathcal C}Y)
    \bigr).
\end{aligned}
\]
Postcomposition with
\[
    \chi_Y
    \colon
    F(D_{\mathcal C}Y)
    \xrightarrow{\;\simeq\;}
    D_{\mathcal D}F(Y)
\]
gives the displayed morphism
\[
    \mathbb B_{\mathcal C}(X,Y)
    \longrightarrow
    \mathbb B_{\mathcal D}
    \bigl(
        F(X),
        F(Y)
    \bigr).
\]

The compatibility of $\chi$ with biduality implies that these maps commute
with transposition. Explicitly, for
\[
    q
    \colon
    X
    \longrightarrow
    D_{\mathcal C}Y,
\]
the transposition of its image is canonically equivalent to the image of
the transposition
\[
    D_{\mathcal C}(q)
    \circ
    \operatorname{bid}_Y
    \colon
    Y
    \longrightarrow
    D_{\mathcal C}X.
\]
The full coherent compatibility follows from the coherent
duality-preserving structure on $F$.

Hence the diagonal map
\[
    \mathbb B_{\mathcal C}(X,X)
    \longrightarrow
    \mathbb B_{\mathcal D}
    \bigl(
        F(X),
        F(X)
    \bigr)
\]
is $C_2$-equivariant. Passing to homotopy fixed points gives a natural
transformation
\[
    \mathcal Q_{\mathcal C}^{\mathrm s}(X)
    \longrightarrow
    \mathcal Q_{\mathcal D}^{\mathrm s}(F(X)).
\]
Naturality in $X$ gives
\[
    \mathcal Q_{\mathcal C}^{\mathrm s}
    \longrightarrow
    \mathcal Q_{\mathcal D}^{\mathrm s}
    \circ
    F^{\op}.
\]

By unstraightening, this natural transformation induces a functor between
the right fibrations of forms
\[
    \operatorname{He}
    \bigl(
        \mathcal C,
        \mathcal Q_{\mathcal C}^{\mathrm s}
    \bigr)
    \longrightarrow
    \operatorname{He}
    \bigl(
        \mathcal D,
        \mathcal Q_{\mathcal D}^{\mathrm s}
    \bigr).
\]
Passing to maximal $\infty$-groupoids gives the functor on spaces of
forms.

If
\[
    q^\sharp
    \colon
    X
    \xrightarrow{\;\simeq\;}
    D_{\mathcal C}X
\]
is an equivalence, then its transported adjoint is
\[
    F(X)
    \xrightarrow{F(q^\sharp)}
    F(D_{\mathcal C}X)
    \xrightarrow{\chi_X}
    D_{\mathcal D}F(X).
\]
Both factors are equivalences. Hence the transported form is
nondegenerate, and the functor on spaces of forms restricts to the stated
functor on spaces of Poincaré objects.

The identity and composition coherences follow from the corresponding
coherences of the duality comparisons.
\end{proof}


\begin{proposition}[Pullback of Hermitian and Poincaré perfect objects]
\label{prop:pullback-hermitian-perfect-objects}
Let
\[
    p
    \colon
    X
    \longrightarrow
    Y
\]
be a morphism in
\[
    \PSh(\Aff).
\]

Pullback induces a natural transformation
\[
    \mathcal Q_Y^{\mathrm s}
    \longrightarrow
    \mathcal Q_X^{\mathrm s}
    \circ
    (p^*)^{\op}
\]
and consequently a functor
\[
    p^*
    \colon
    \operatorname{Pn}
    \bigl(
        \Perf(Y),
        \mathcal Q_Y^{\mathrm s}
    \bigr)
    \longrightarrow
    \operatorname{Pn}
    \bigl(
        \Perf(X),
        \mathcal Q_X^{\mathrm s}
    \bigr).
\]

For a Hermitian perfect object
\[
    (M,h_M)
\]
on $Y$, the underlying Hermitian equivalence on its pullback is
\[
\begin{aligned}
    h_{p^*M}
    \colon
    p^*M
    &\xrightarrow{p^*(h_M)}
    p^*(M^\vee)
    \\
    &\xrightarrow{\chi_{p,M}}
    \bigl(p^*M\bigr)^\vee,
\end{aligned}
\]
where
\[
    \chi_{p,M}
    \colon
    p^*(M^\vee)
    \xrightarrow{\;\simeq\;}
    \bigl(p^*M\bigr)^\vee
\]
is the canonical duality comparison.

Under the equivalences
\[
    \operatorname{Pn}
    \bigl(
        \Perf(Y),
        \mathcal Q_Y^{\mathrm s}
    \bigr)
    \simeq
    \bigl(
        \Perf(Y)^\simeq
    \bigr)^{hC_2}
\]
and
\[
    \operatorname{Pn}
    \bigl(
        \Perf(X),
        \mathcal Q_X^{\mathrm s}
    \bigr)
    \simeq
    \bigl(
        \Perf(X)^\simeq
    \bigr)^{hC_2},
\]
this functor identifies with the functor on homotopy fixed points induced
by
\[
    p^*
    \colon
    \Perf(Y)^\simeq
    \longrightarrow
    \Perf(X)^\simeq.
\]
\end{proposition}

\begin{proof}
By
\Ref{cor:pullback-perfect-duality},
the functor
\[
    p^*
    \colon
    \Perf(Y)
    \longrightarrow
    \Perf(X)
\]
is an exact duality-preserving strong symmetric monoidal functor, with
duality comparison
\[
    \chi_{p,M}
    \colon
    p^*(M^\vee)
    \xrightarrow{\;\simeq\;}
    \bigl(p^*M\bigr)^\vee.
\]

Apply
\Ref{prop:functoriality-homotopy-symmetric-structures}.
It gives the natural transformation
\[
    \mathcal Q_Y^{\mathrm s}
    \longrightarrow
    \mathcal Q_X^{\mathrm s}
    \circ
    (p^*)^{\op}
\]
and the induced functor on spaces of Poincaré objects.

For a Hermitian equivalence
\[
    h_M
    \colon
    M
    \xrightarrow{\;\simeq\;}
    M^\vee,
\]
the transported adjoint is
\[
    p^*M
    \xrightarrow{p^*(h_M)}
    p^*(M^\vee)
    \xrightarrow{\chi_{p,M}}
    \bigl(p^*M\bigr)^\vee.
\]
The coherent duality-preserving structure on $p^*$ transports the complete
homotopy-fixed-point datum, so this composite is the underlying map of a
Hermitian structure on $p^*M$.

The fixed-point description follows from the naturality of the
equivalence of
\Ref{prop:hermitian-as-poincare}
with respect to duality-preserving exact functors.
\end{proof}


\begin{proposition}[Affine base change of Hermitian and Poincaré modules]
\label{prop:base-change-hermitian-perfect-modules}
Let
\[
    f
    \colon
    A
    \longrightarrow
    A'
\]
be a morphism in
\[
    \CAlg(\V),
\]
and let
\[
    F
    :=
    -\otimes_A A'
    \colon
    \Perf(A)
    \longrightarrow
    \Perf(A').
\]

No equivariance condition on the commutative base is required.

Extension of scalars induces a natural transformation
\[
    \mathcal Q_A^{\mathrm s}
    \longrightarrow
    \mathcal Q_{A'}^{\mathrm s}
    \circ
    F^{\op}
\]
and a functor
\[
    F
    \colon
    \operatorname{Pn}
    \bigl(
        \Perf(A),
        \mathcal Q_A^{\mathrm s}
    \bigr)
    \longrightarrow
    \operatorname{Pn}
    \bigl(
        \Perf(A'),
        \mathcal Q_{A'}^{\mathrm s}
    \bigr).
\]

Write
\[
    \chi_M
    \colon
    F(M^\vee)
    \xrightarrow{\;\simeq\;}
    F(M)^\vee
\]
for the canonical duality comparison. If
\[
    h_M
    \colon
    M
    \xrightarrow{\;\simeq\;}
    M^\vee
\]
is a Hermitian structure, then the underlying Hermitian equivalence on
the base-changed module is
\[
    h_{F(M)}
    :=
    \chi_M
    \circ
    F(h_M)
    \colon
    F(M)
    \xrightarrow{\;\simeq\;}
    F(M)^\vee.
\]
\end{proposition}

\begin{proof}
By
\Ref{cor:base-change-perfect-duality},
extension of scalars
\[
    F
    =
    -\otimes_A A'
    \colon
    \Perf(A)
    \longrightarrow
    \Perf(A')
\]
is an exact duality-preserving strong symmetric monoidal functor.

Apply
\Ref{prop:functoriality-homotopy-symmetric-structures}.
It supplies the natural transformation
\[
    \mathcal Q_A^{\mathrm s}
    \longrightarrow
    \mathcal Q_{A'}^{\mathrm s}
    \circ
    F^{\op}
\]
and the induced functor on spaces of Poincaré objects.

For a Hermitian equivalence
\[
    h_M
    \colon
    M
    \xrightarrow{\;\simeq\;}
    M^\vee,
\]
the transported adjoint is
\[
    F(M)
    \xrightarrow{F(h_M)}
    F(M^\vee)
    \xrightarrow{\chi_M}
    F(M)^\vee.
\]
The coherent compatibility of $\chi$ with biduality transports the entire
homotopy-fixed-point datum of $h_M$. Hence this composite defines a
Hermitian structure on
\[
    F(M)
    =
    M\otimes_A A'.
\]
\end{proof}
\section{Alternative refinements and examples}
\label{sec:examples-perfect-poincare}

\subsection{Homotopy-symmetric and homotopy-quadratic refinements}

\begin{remark}
\label{rem:poincare-infinity-category}
A stable category with coherent duality does not determine a unique
Poincar\'e refinement. It canonically determines at least two
distinguished refinements associated to the symmetric bilinear functor
\[
    \mathbb B_A:
\]
\[
    \mathcal Q_A^{\mathrm s}(M)
    =
    \mathbb B_A(M,M)^{hC_2},
\]
the homotopy-symmetric refinement, and
\[
    \mathcal Q_A^{\mathrm q}(M)
    =
    \mathbb B_A(M,M)_{hC_2},
\]
the homotopy-quadratic refinement.

Homotopy orbits are colimit-preserving and therefore exact on spectra.
They preserve finite cartesian cubes. Hence
\[
    \mathcal Q_A^{\mathrm q}
\]
is again reduced and
\[
    2
\]
-excisive.

Both refinements have symmetric bilinear cross-effect equivalent to
\[
    \mathbb B_A
\]
and therefore represent the same underlying coherent duality. Further
intermediate, genuine, or enhanced refinements may require additional
structure. See
\cite{CalmesDottoHarpazEtAl}.
\end{remark}


\subsection{Examples}

\begin{example}[Perfect complexes]
\label{ex:perfect-complexes}
Let
\[
    R
\]
be an $\EE_\infty$-ring spectrum, including the case
\[
    R=HR_0
\]
for an ordinary commutative ring
\[
    R_0,
\]
and take
\[
    \V
    =
    \Mod_R(\Sp).
\]
For a commutative
\[
    R
\]
-algebra
\[
    A,
\]
the category
\[
    \Perf(A)
\]
is the rigid stable category of perfect
\[
    A
\]
-modules. The canonical duality is
\[
    M^\vee
    =
    \HHom_A(M,A).
\]
A Hermitian perfect module in the sense of this chapter is a perfect
module equipped with a coherent nondegenerate symmetric equivalence
\[
    M
    \simeq
    M^\vee.
\]
\end{example}

\begin{example}[Complex perfect complexes]
\label{ex:complex-perfect-complexes}
Let
\[
    \V
    =
    \Mod_{H\mathbb C}(\Sp).
\]
Then
\[
    \Perf(H\mathbb C)
\]
is the stable $\infty$-category of perfect complexes of complex vector
spaces, and
\[
    \HHom_{\mathbb C}(M,N)
    \simeq
    M^\vee\otimes_{\mathbb C}N.
\]
The Hermitian forms defined here are nondegenerate symmetric
complex-bilinear self-dualities. They are not sesquilinear forms, since no
complex conjugation has been introduced.
\end{example}

The rigid theory developed here is internal to symmetric monoidal module
categories. The next chapter passes from individual perfect modules to
associative algebra objects and from morphisms to bimodules. In that
Morita-theoretic setting, dualizability is replaced by adjointability, and
the canonical rigid duality is replaced by a coherent duality on a
right-perfect Morita
\[
    (\infty,2)
\]
